%
% 6.006 problem set 5 solutions template
%
\documentclass[12pt,twoside]{article}

\input{macros-sp20}
\newcommand{\theproblemsetnum}{5}

\title{6.006 Problem Set \theproblemsetnum}

\begin{document}

\handout{Problem Set \theproblemsetnum}

\setlength{\parindent}{0pt}
\medskip\hrulefill\medskip

{\bf Name:} BUAA-TYZ

\medskip

{\bf Collaborators:}

\medskip\hrulefill

%%%%%%%%%%%%%%%%%%%%%%%%%%%%%%%%%%%%%%%%%%%%%%%%%%%%%
% See below for common and useful latex constructs. %
%%%%%%%%%%%%%%%%%%%%%%%%%%%%%%%%%%%%%%%%%%%%%%%%%%%%%

% Some useful commands:
%$f(x) = \Theta(x)$
%$T(x, y) \leq \log(x) + 2^y + \binom{2n}{n}$
% {\tt code\_function}


% You can create unnumbered lists as follows:
%\begin{itemize}
%    \item First item in a list
%        \begin{itemize}
%            \item First item in a list
%                \begin{itemize}
%                    \item First item in a list
%                    \item Second item in a list
%                \end{itemize}
%            \item Second item in a list
%        \end{itemize}
%    \item Second item in a list
%\end{itemize}

% You can create numbered lists as follows:
%\begin{enumerate}
%    \item First item in a list
%    \item Second item in a list
%    \item Third item in a list
%\end{enumerate}

% You can write aligned equations as follows:
%\begin{align}
%    \begin{split}
%        (x+y)^3 &= (x+y)^2(x+y) \\
%                &= (x^2+2xy+y^2)(x+y) \\
%                &= (x^3+2x^2y+xy^2) + (x^2y+2xy^2+y^3) \\
%                &= x^3+3x^2y+3xy^2+y^3
%    \end{split}
%\end{align}

% You can create grids/matrices as follows:
%\begin{align}
%    A =
%    \begin{bmatrix}
%        A_{11} & A_{21} \\
%        A_{21} & A_{22}
%    \end{bmatrix}
%\end{align}

% You can include images and PDFs as follows:
% \includegraphics[width=0.5\textwidth]{img.jpg}

\begin{problems}

\problem  % Problem 1

\begin{problemparts}
\problempart % Problem 1a
pass
\problempart % Problem 1b
pass
\problempart % Problem 1c

  \begin{tikzpicture}
  \tikzset{level distance=40pt, sibling distance=30pt}    
  \Tree
  [.\node[draw, circle]{$A$}; 
  [.\node[draw, circle]{$B$}; 
  [.\node[draw, circle]{$C$}; [.\node[draw, circle]{$E$}; ] [.\node[draw, circle]{$F$}; ] ] 
  [.\node[draw, circle]{$D$}; ] 
  ]
  ]
  \end{tikzpicture}
BFS: A B C D E F
  \begin{tikzpicture}
  \tikzset{level distance=40pt, sibling distance=30pt}    
  \Tree
  [.\node[draw, circle]{$A$}; 
  [.\node[draw, circle]{$B$}; 
  [.\node[draw, circle]{$C$}; [.\node[draw, circle]{$E$}; ] 
  [.\node[draw, circle]{$F$}; [.\node[draw, circle]{$D$}; ] ] ] 
  ]
  ]
  \end{tikzpicture}
 DFS: A B C E F D
\problempart % Problem 1d
\begin{enumerate}
  \item Remove D->F: {A, B, C, {E, F}, D}
  \item Remove F->D: {A, B, C, D, E, F}
\end{enumerate}
\end{problemparts}

\newpage
\problem  % Problem 2
For each plant, it builds an undirected graph with supported buildings. 
Do a BFS for each plant, which costs $O(|V|+|E|)$. And $|E|<=\binom{n^2}{2}=O(n^4)$,
we salely conclude that this algorithm safisfies the need.

\problem  % Problem 3
If we think about the conception: Indegree, then it's easy. For a robot, if it's Indegree
is two, then it's impossible to achieve the goal. So we just calculate the Indegree of each
robot and judge whether it is smaller than 2. The algorithm scans all pairs once, which takes $O(n)$

\problem  % Problem 4
pass

\problem  % Problem 5


\problem  % Problem 6

\begin{problemparts}
\problempart % Problem 6a
\problempart % Problem 6b
\problempart % Problem 6c
\problempart Submit your implementation to {\small\url{alg.mit.edu}}.
\end{problemparts}

\end{problems}

\end{document}
